\section{Trust Region \ac*{MPC}}
\label{sec:tr_mpc}

Inspired by the success of trust region methods in \ac{RL}, we propose a trust region formulation for sampling-based \ac{MPC} to improve the optimization process and remove the need for heuristic parameter adaptation.

At iteration $j$, let $\pdf_j$ denote the current proposal \ac{PDF} over control sequences $\vVarCem$, and $\pdf_{j+1}$ the next proposal \ac{PDF} to be optimized.
The trust region optimization problem is then given by
\begin{align}
    \label{eq:tr_objective_function}
    &\min_{\pdf_{j+1}} \Eop_{\pdf_{j+1}}\{J(\vVarCem)\}
    % \\
    % &
    \quad \text{s.t.} \quad D_{\mathrm{KL}}(\pdf_{j+1} \| \pdf_j) \leq \epsilon \enspace, 
    \\
    &
    \quad H(\pdf_{j+1}) \ge H_{\min} \enspace,
    % \\
    % &
    \quad  \Eop_{\pdf_{j+1}} \{ 1\} = 1 \enspace,
\end{align}
where $D_{\mathrm{KL}}(\cdot \| \cdot)$ is the \ac{KL} divergence measuring the distributional change, $H(\cdot)$ is the entropy, $\epsilon>0$ is the maximum allowed \ac{KL} divergence between consecutive proposals,  $H_{\min}$ is the minimum required entropy of the proposal, and the last constraint ensures that $\pdf_{j+1}$ is normalized.
The \ac{KL} constraint prevents overly aggressive updates, while the minimum entropy constraint encourages exploration and mitigates premature collapse of the proposal.
Note that the entropy constraint is optional and can be omitted if task-specific knowledge suggests that proposal collapse is not a concern, \eg, when a small number of iterations is used.

To solve the constrained problem, we form the Lagrangian 
\begin{align}
    & \mathcal{L}(\pdf_{j+1},\eta,\alpha, \nu) = \Eop_{\pdf_{j+1}}\{J(\vVarCem)\}
    % \\ 
    % & \quad 
    + \eta\left(D_{\mathrm{KL}}(\pdf_{j+1} \| \pdf_j) - \epsilon\right) 
    \\ & \quad
    \label{eq:tr:lagrangian}
    + \alpha\left(H_{\min} - H(\pdf_{j+1})\right) 
    + \nu \left( \Eop_{\pdf_{j+1}} \{ 1\} - 1 \right)
    \enspace,
\end{align}
where $\eta$, $\alpha$, and $\nu$ are the Lagrange multipliers associated with the \ac{KL} divergence, entropy, and normalization constraints, respectively.
To find the optimal proposal for the next iteration, we take the functional derivative of $\mathcal{L}$ w.r.t. $\pdf_{j+1}$ and set it to zero (see \cref{sec:appendix:lagrangian} for details).
% The normalization constraint introduces only an additive constant in the log-domain; consequently, its associated multiplier $\nu$ is implicitly handled and omitted from the subsequent derivation steps (see \cref{sec:appendix:lagrangian} for details).
The optimal proposal is then given by
\begin{align}
    \label{eq:opt_proposal}
    \pdf_{j+1}^{\ast}(\vVarCem) &= \frac{1}{Z_{\eta,\alpha}}\; \pdf_j(\vVarCem)^{\frac{\eta}{\eta+\alpha}}\exp\left(-\frac{J(\vVarCem)}{\eta+\alpha}\right)
    \enspace,
    \\
    \label{eq:opt_proposal:normalization}
    Z_{\eta,\alpha} 
    % &= \int \pdf_j(\vVarCem)^{\tfrac{\eta}{\eta+\alpha}}\exp\left(-\frac{J(\vVarCem)}{\eta+\alpha}\right)d\vVarCem
    &= \Eop_{\pdf_j}\left\{ \pdf_j(\vVarCem)^{\tfrac{\eta}{\eta+\alpha}-1}\exp \left( -\frac{J(\vVarCem)}{\eta+\alpha} \right) \right\} \enspace,
\end{align}
where $Z_{\eta,\alpha}$ is the normalization constant, written as expectation over $\pdf_j$.
% (we will use this to establish a sample-based estimate of $Z_{\eta,\alpha}$).
Since the multiplier $\nu$ is implicitly handled in the normalization, it can be omitted from subsequent derivation steps.
Note that in the case of a Gaussian proposal $\pdf_j(\vVarCem)$, for $\alpha>0$, $\pdf_{j+1}^{\ast}$ is generally non-Gaussian.
However, for $\alpha=0$, and convex quadratic $J(\vVarCem)$, Gaussianity is preserved.

By substituting $\pdf_{j+1}^{\ast}$ back into the simplified Lagrangian $\mathcal{L}(\pdf_{j+1}^{\ast},\eta,\alpha)$ (see \cref{sec:appendix:dual}), we obtain the dual function
\begingroup
% \medmuskip=3.mu% default is 4mu plus 2mu minus 4mu
% \thickmuskip=3.mu% default is 5mu plus 5mu
% \thinmuskip=2.mu%default is 3mu
\begin{align}
    \label{eq:tr_dual}
    g(\eta, \alpha) 
    % =  \mathcal{L}(\pdf_{j+1}^{\ast},\eta,\alpha) 
    = -\eta \epsilon + \alpha H_{\min} - (\eta + \alpha) \log Z_{\eta,\alpha} \enspace,
\end{align}
\endgroup
which we use to find the optimal Lagrange multipliers $\eta^{\ast}$ and $\alpha^{\ast}$ by numerically solving the dual problem
\begin{align}
    \label{eq:tr_dual:opt_problem}
    \max_{\eta \geq 0, ~\alpha \geq 0} g(\eta, \alpha) \enspace.
\end{align}
Note that, in our case, strong duality holds, so that solving the dual problem is equivalent to solving the primal problem~\cref{eq:tr_objective_function} (see \cref{sec:appendix:strong_duality}) and that $\pdf_{j+1}$ never degenerates, \ie, $\eta^{\ast} + \alpha^{\ast} > 0$ (see \cref{sec:appendix:pos_multiplier_sum}).

Since the optimization over $\eta$ and $\alpha$ is always two-dimensional (in particular, it is independent of the control input dimension and horizon length), standard gradient-based solvers are efficient in practice.
% The gradient of the dual function $\nabla g(\eta,\alpha) = [ \nicefrac{\partial g}{\partial \eta}, \nicefrac{\partial g}{\partial \alpha} ]\T$ is given by
% \begin{align}
%     \label{eq:tr_dual:gradient}
%     \begin{bmatrix}
%         -\log Z_{\eta,\alpha}-\frac{\alpha\,\Eop_{\pdf_{j+1}^{\ast}}[\log \pdf_j(\vVarCem)] + \Eop_{\pdf_{j+1}^{\ast}}[J(\vVarCem)]}{\eta + \alpha}-\epsilon 
%         \\
%         -\log Z_{\eta,\alpha}+\frac{\eta\,\Eop_{\pdf_{j+1}^{\ast}}[\log \pdf_j(\vVarCem)] - \Eop_{\pdf_{j+1}^{\ast}}[J(\vVarCem)]}{\eta + \alpha}+H_{\min}
%     \end{bmatrix}
%     \enspace.
% \end{align}
% Note that by substituting $\pdf_{j+1}^{\ast}$ from~\cref{eq:opt_proposal} and $Z_{\eta,\alpha}$ from \cref{eq:opt_proposal:normalization} into~\cref{eq:tr_dual:gradient}, the gradient can be entirely written in terms of expectations of $\pdf_j$.
% Consequently, the gradient~\cref{eq:tr_dual:gradient} can be estimated from samples drawn from $\pdf_j$.
The gradient of the dual function $\nabla g(\eta,\alpha)$ is given by
\begin{align}
    \label{eq:tr_dual:gradient}
    \begin{bmatrix}
        \frac{\partial g}{\partial \eta}
        \\
        \frac{\partial g}{\partial \alpha}
    \end{bmatrix}
    =
    \begin{bmatrix}
        -\log Z_{\eta,\alpha}
        -(\eta+\alpha)\frac{\partial \log Z_{\eta,\alpha}}{\partial \eta}
        -\epsilon
        \\
        -\log Z_{\eta,\alpha}
        -(\eta+\alpha)\frac{\partial \log Z_{\eta,\alpha}}{\partial \alpha}
        +H_{\min}
    \end{bmatrix}
\end{align}
where the partial derivatives of $\log Z_{\eta,\alpha}$ are
\begin{align}
    &
    \frac{\partial \log Z_{\eta,\alpha}}{\partial \eta}
    =
    \frac{1}{(\eta+\alpha)^2}
    \\
    &
    \enspace\cdot
    \frac{\Eop_{\pdf_{j}} \left\{
    \pdf_{j}(\vVarCem)^{\frac{\eta}{\eta+\alpha}-1}
    \exp \left(-\frac{J(\vVarCem)}{\eta+\alpha}\right)
    \left(\alpha\log \pdf_{j}(\vVarCem)+J(\vVarCem)\right)
    \right\} }
    {
        % \Eop_{\pdf_{j}} \left\{
        % \pdf_{j}(x)^{\frac{\eta}{\eta+\alpha}-1}
        % \exp \left(-\frac{J(x)}{\eta+\alpha}\right)
        % \right\} 
        Z_{\eta,\alpha}
    } \, ,
    % =\frac{1}{(\eta+\alpha)^2}
    % \Eop_{\pdf_{j+1}^{\ast}} \left[\alpha\log \pdf_{j}(x)+J(x)\right].
    \\
    &
    \frac{\partial \log Z_{\eta,\alpha}}{\partial \alpha}
    =
    \frac{1}{(\eta+\alpha)^2}
    \\
    &
    \enspace\cdot
    \frac{\Eop_{\pdf_{j}} \left\{
    \pdf_{j}(\vVarCem)^{\frac{\eta}{\eta+\alpha}-1}
    \exp \left(-\frac{J(\vVarCem)}{\eta+\alpha}\right)
    \left(-\eta\log \pdf_{j}(\vVarCem)+J(\vVarCem)\right)
    \right\} }
    {
        % \Eop_{\pdf_{j}} \left\{
        % \pdf_{j}(x)^{\frac{\eta}{\eta+\alpha}-1}
        % \exp \left(-\frac{J(x)}{\eta+\alpha}\right)
        % \right\} 
        Z_{\eta,\alpha}
    } \, .
    % =\frac{1}{(\eta+\alpha)^2}
    % \Eop_{\pdf_{j+1}^{\ast}} \left[-\eta\log \pdf_{j}(x)+J(x)\right].
\end{align}
Note that the gradient is entirely written in terms of expectations of $\pdf_j$.
Consequently, the gradient can be estimated from samples drawn from $\pdf_j$.

% Once the optimal Lagrange multipliers $\eta^{\ast}$ and $\alpha^{\ast}$ have been obtained, the optimal proposal \ac{PDF} $\pdf_{j+1}^{\ast}$ follows directly from~\cref{eq:opt_proposal}.
% We can derive a sampling-based approximation of $\pdf_{j+1}^{\ast}$ based on samples drawn from $\pdf_j$ and weights
% \begin{align}
%     \label{eq:tr_dual:weights}
%     w_i \propto \pdf_j(\vVarCem^{(i)})^{\frac{\eta^{\ast}}{\eta^{\ast}+\alpha^{\ast}}-1}\exp\left(-\frac{J(\vVarCem^{(i)})}{\eta^{\ast}+\alpha^{\ast}}\right) \enspace.
% \end{align}
% Although different in derivation, the resulting weights have a similar form as the update in \ac{MPPI}-like methods, but with a principled weighting scheme derived from the trust region formulation, and without the need for heuristic temperature adaptation. 
% When the entropy constraint is inactive, \ie, $\alpha=0$, the weights reduce to exactly the same form as in \cref{eq:imp_weights}.

% Note that for $\alpha>0$, the \ac{PDF} $\pdf_{j+1}^{\ast}$ is non-Gaussian even if the previous proposal $\pdf_j$ is Gaussian, due to the exponentiation of $\pdf_j$.
% Nevertheless, the moment-matching approximation projects $\pdf_{j+1}^{\ast}$ back onto the Gaussian family.


Once the optimal Lagrange multipliers $\eta^{\ast}$ and $\alpha^{\ast}$ have been obtained, the optimal proposal \ac{PDF} $\pdf_{j+1}^{\ast}$ follows directly from~\cref{eq:opt_proposal}.
We can derive a Dirac-mixture approximation of $\pdf_{j+1}^{\ast}$ based on samples drawn from $\pdf_j$ and weights
\begin{align}
    \label{eq:tr_dual:weights}
    w_i \propto \pdf_j(\vVarCem^{(i)})^{\frac{\eta^{\ast}}{\eta^{\ast}+\alpha^{\ast}}-1}\exp\left(-\frac{J(\vVarCem^{(i)})}{\eta^{\ast}+\alpha^{\ast}}\right) \enspace.
\end{align}
Finally, we project our Dirac-mixture approximation of $\pdf_{j+1}^{\ast}$ onto the same parametric form as the original proposal \ac{PDF}, \eg, via moment matching for Gaussians, or expectation--maximization for Gaussian mixtures.
In this work, we use Gaussian proposals, which yield a closed-form update for the new mean $\evVarCem_{j+1}$ and covariance $\mC_{j+1}$, expressed in terms of the weighted empirical sample mean and covariance.
Notably, when $\alpha^{\ast} = 0$, this projection does not introduce additional errors, as the true $\pdf_{j+1}^{\ast}$ is Gaussian. 
In the general case, however, the projection incurs an additional approximation error.

Although different in derivation, the resulting weights \cref{eq:tr_dual:weights} exhibit a form similar to those in \cref{eq:imp_weights} used by \ac{MPPI}-like methods.
In particular, when the entropy constraint is inactive, \ie, $\alpha^{\ast}=0$, the two formulations coincide exactly.
Importantly, however, our approach yields optimal weights in a principled manner from the trust region formulation, thereby eliminating the need for heuristic hyperparameter adaptation.